\documentclass[../../main.tex]{subfiles}
\graphicspath{{../../image/}{../../figure/}}
\begin{document}

\section{标量守恒律}\label{section:ch17-标量守恒律}

这里将看到上一节勾勒的纲领如何在标量守恒律
\begin{align}
\partial_t u+\partial_x f(u)=0
\label{align:ch17-1741}
\end{align}
的情形下实现。

\begin{theorem}[标量情形的补偿紧致性]\label{theorem:ch17-标量补偿紧致}
设 $\Omega$ 是 $\mathbb{R}^2$ 的开子集，$\{u_k(x,t)\}$ 是 $L^\infty(\Omega)$ 中的有界序列，满足
\begin{align}
\partial_t\eta(u_k)+\partial_x q(u_k)\subset\text{紧集于 }W^{-1,2}_{\mathrm{loc}}(\Omega),
\label{align:ch17-1742}
\end{align}
其中 $(\eta,q)$ 是\eqref{align:ch17-1741}的任意熵--熵流对。则存在子列 $\{u_j\}$，使得
\begin{align}
u_j\rightharpoonup\bar u,\qquad f(u_j)\rightharpoonup f(\bar u),\qquad j\to\infty,
\label{align:ch17-1743}
\end{align}
在 $L^\infty$ 弱$^*$ 意义下成立。此外，若使 $f''(u)\neq0$ 的 $u$ 之集在 $\mathbb{R}$ 中稠密，则 $\{u_j\}$ 在 $\Omega$ 上几乎处处收敛于 $\bar u$。
\end{theorem}

\begin{proof}
应用\cref{theorem:ch17-young测度存在性}，提取子列 $\{u_j\}$ 并联系 Young 测度族 $\nu=\nu_{x,t}$，使得对任意连续函数 $h$，$h(u_j)\rightharpoonup\langle\nu,h\rangle$。于是 $u_j\rightharpoonup\bar u=\langle\nu,u\rangle$，$f(u_j)\rightharpoonup\langle\nu,f\rangle$。我们只需证明 $\langle\nu,f\rangle=f(\bar u)$，并且当不存在使 $f'(u)$ 为常数的区间时，$\nu$ 退化为 Dirac 质量。

我们对特殊的熵--熵流对 $(u,f(u))$ 与 $(f(u),g(u))$ 使用\eqref{align:ch17-17310}，其中
\begin{align}
g(u)=\int_0^u[f'(v)]^2\,dv,
\label{align:ch17-1744}
\end{align}
从而得到
\begin{align}
\langle\nu,u\rangle\langle\nu,g\rangle-\langle\nu,f\rangle\langle\nu,f\rangle=\langle\nu,ug-f^2\rangle.
\label{align:ch17-1745}
\end{align}

由 Schwarz 不等式，
\begin{align}
[f(u)-f(\bar u)]^2\leqslant(u-\bar u)[g(u)-g(\bar u)],
\label{align:ch17-1746}
\end{align}
我们推出
\begin{align}
\bigl\langle\nu,[f(u)-f(\bar u)]^2-(u-\bar u)[g(u)-g(\bar u)]\bigr\rangle\leqslant0.
\label{align:ch17-1747}
\end{align}
利用\eqref{align:ch17-1745}，\eqref{align:ch17-1747}化为% 原书此处误排为 (16.4.5),(16.4.7)
\begin{align}
\bigl[\langle\nu,f\rangle-f(\bar u)\bigr]^2\leqslant0,
\label{align:ch17-1748}
\end{align}
由此 $\langle\nu,f\rangle=f(\bar u)$。特别地，\eqref{align:ch17-1747}左端将为零。于是\eqref{align:ch17-1746}在 $\nu$ 的支撑上对 $u$ 必以等式成立。然而，Schwarz 不等式\eqref{align:ch17-1746}仅当 $f'$ 在以 $\bar u$ 与 $u$ 为端点的区间上为常数时才以等式成立。当不存在这样的区间时，$\nu$ 的支撑坍缩为单点，$\nu$ 退化为 Dirac 质量 $\delta_{\bar u}$。证明完成。
\end{proof}

如前一节所指出的，可由消失黏性方法生成一个满足\cref{theorem:ch17-标量补偿紧致}假设的序列 $\{u_k\}$：令 $u_k=u_{\mu_k}$，$\mu_k\to0$（$k\to\infty$），其中 $u_\mu$ 是
\begin{align}
\partial_t u+\partial_x f(u)=\mu\,\partial_x^2u
\label{align:ch17-1749}
\end{align}
在上半平面上以
\begin{align}
u(x,0)=u_{0\mu}(x),\qquad-\infty<x<\infty,
\label{align:ch17-17410}
\end{align}
为初值、且 $u_{0\mu}$ 在 $L^\infty(-\infty,\infty)\cap L^2(-\infty,\infty)$ 中一致有界的解。事实上，所得 $\{u_k\}$ 在 $L^\infty$ 中有界，因为由最大值原理 $\|u_\mu\|_{L^\infty}\leqslant\|u_{0\mu}\|_{L^\infty}$。此外，由17.3节的一般论证，\eqref{align:ch17-1742}对所有熵--熵流对 $(\eta,q)$ 成立；该处论证在此特别适用，因为\eqref{align:ch17-1741}配备了一致凸熵 $u^2$。最后，当 $\mu\downarrow0$ 时，$\mu\,\partial_x^2u_\mu\to0$ 在分布意义下成立。于是得到下述

\begin{theorem}[标量守恒律的消失黏性极限]\label{theorem:ch17-标量消失黏性}
设当 $\mu\downarrow0$ 时，$u_{0\mu}\rightharpoonup u_0$ 在 $L^\infty(-\infty,\infty)$ 弱$^*$ 意义下成立。则存在序列 $\{\mu_j\}$（$\mu_j\to0$，$j\to\infty$），使得\eqref{align:ch17-1749}、\eqref{align:ch17-17410}的解序列 $\{u_{\mu_j}\}$ 在 $L^\infty$ 弱$^*$ 意义下收敛于某函数 $\bar u$，它是上半平面上\eqref{align:ch17-1741}的解，且初值为 $\bar u(x,0)=u_0(x)$（$-\infty<x<\infty$）。此外，若使 $f''(u)\neq0$ 的 $u$ 之集在 $\mathbb{R}$ 中稠密，则 $\{u_{\mu_j}\}$ 或其子列在上半平面上几乎处处收敛于 $\bar u$。
\end{theorem}

\end{document}
